\section{The complete strip Green potential}\label{sec:green}

\subsection{The finite-factor identity}

Fix
\begin{equation}\label{eq:kernels}
 a=\frac\pi8,\qquad K(t)=\frac1{16\cosh(at)},\qquad
 G(t,y)=\atanh\frac{\sin(ay)}{\cosh(at)}.
\end{equation}
Each copy of \(K\) has mass \(1/2\). Let \(\mu_X\) be the probability
measure on the two lines \(y=0,8\) that has density \(K(x-X)\) on
each. For an admissible function \(f\), write
\[
 \calJ[f](X)=\int\log|f|\dd\mu_X-\log|f(X+4\ii)|,
 \qquad J_\lambda(X)=\calJ[H_\lambda](X).
\]

The Green function of \(0<\Im z<B\), normalized by
\(-\Delta G_B=2\pi\delta\), is
\begin{equation}\label{eq:strip-green}
 G_B(z,\zeta)=\frac12\log
 \frac{\cosh(\pi(\Re z-\Re\zeta)/B)-\cos(\pi(\Im z+\Im\zeta)/B)}
      {\cosh(\pi(\Re z-\Re\zeta)/B)-\cos(\pi(\Im z-\Im\zeta)/B)}.
\end{equation}
Indeed \(z\mapsto\ee^{\pi z/B}\) maps the strip to the upper
half-plane; the logarithm of the reflected-distance ratio gives
\eqref{eq:strip-green}. Its singularity is \(-\log|z-\zeta|\).
At \(B=8\), observation \(X+4\ii\), its value is
\(G(X-\Re\zeta,\Im\zeta)\). Boundary normal derivatives divided
by \(2\pi\) give exactly the two densities \(K\).

For an interior zero \(\zeta\), the singularities of
\(\log|z-\zeta|+G_B(z,\zeta)\) cancel. This harmonic function has the
linear factor's boundary logarithm. Truncating the strip vertically
proves its Poisson representation: the side harmonic measures decay
exponentially while these logarithms grow only logarithmically.
It follows that \(\calJ[z-\zeta]\) is the positive Green value.
An exterior zero gives zero. A boundary zero is handled by its
locally integrable logarithm, or by a limit from the exterior; it
gives no half-weight atom. Finite multiplicities cause no change.

Apply this identity to every complete opposite pair in a finite
truncation of \eqref{eq:product}. Scalar logarithms cancel because
the boundary masses add to one. The only remaining issue is passage
to the complete product.

\subsection{Absolute logarithmic convergence}

Choose
\[
 m_0=\inf_\lambda H_\lambda(0)>0,\qquad
 M_1\ge\sup_{\lambda,|z|\le1}|H_\lambda'(z)|,\quad M_1>0,
\]
and set
\[
 r_*=\min(1,m_0/(2M_1)),\quad
 C_0=\sup_\lambda|\log H_\lambda(0)|,\quad
 C_*=C_0+20C(1+|\log r_*|).
\]
The segment integral of \(H'\) shows that no zero lies in
\(|z|\le r_*\). All constants are uniform for \(0\le\lambda\le1/2\).

\begin{lemma}[Unit-interval logarithmic bound]\label{lem:unit-log}
For \(v\in\R\), \(0\le b\le8\), and \(Q=|v|+12\),
\begin{equation}\label{eq:unit-log}
 \int_v^{v+1}|\log|H_\lambda(x+\ii b)||\dd x
 \le C_*Q\log^2(3Q).
\end{equation}
Moreover the series of paired factor logarithms is absolutely
convergent in local \(L^1\).
\end{lemma}
\begin{proof}
On the interval, \(|z|\le Q\). For a representative \(|\zeta|\le2Q\),
split the logarithm into those of \(z-\zeta\), \(z+\zeta\), and
\(-2\log|\zeta|\). Each positive distance logarithm is at most
\(\log(3Q)\). Its negative part has integral at most two, since
\(|x+\ii b-\zeta|\ge|x-\Re\zeta|\) and
\(\int_{-1}^1-\log|t|\dd t=2\).
The factor therefore costs at most
\(4\log(3Q)+4+2|\log r_*|\).
There are at most \(CQ\log(2Q+2)\) such representatives, by
\eqref{eq:count} and opposite pairing.

For \(|\zeta|>2Q\),
\[
 \left|\log\left|1-\frac{z^2}{\zeta^2}\right|\right|
 \le\frac43\frac{Q^2}{|\zeta|^2}.
\]
Partial summation with half the full count bounds their sum by
\[
 \frac23CQ\{\log(2Q+2)+1\}.
\]
These bounds, together with \(C_0\), give \eqref{eq:unit-log}.
They bound the sum of absolute factor logarithms, which proves the
stronger local \(L^1\) assertion.
\end{proof}

\begin{theorem}[Complete positive potential]\label{thm:green}
For every real \(X\) and \(0\le\lambda\le1/2\), both boundary
logarithms are absolutely integrable against \(K(x-X)\), and
\begin{equation}\label{eq:green-identity}
 J_\lambda(X)=
 \int_\R K(x-X)\{\log|H_\lambda(x)|+\log|H_\lambda(x+8\ii)|\}\dd x
 -\log|H_\lambda(X+4\ii)|
 =\sum_{j:y_j>0}G(X-x_j,y_j).
\end{equation}
The upper-zero sum is absolutely convergent, with all multiplicities.
There is no omitted harmonic term. The potential is nonnegative and
unchanged by multiplication of the field by a nonzero scalar.
\end{theorem}
\begin{proof}
The polynomial-logarithmic bound \eqref{eq:unit-log} is summable
against the exponentially decaying Poisson weights. Thus the finite
paired-product identity passes to the boundary integrals in \(L^1\).
At \(X+4\ii\), the product converges away from every zero by
\eqref{eq:zero-strip}. This proves the complete identity. Its finite
Green sums are nonnegative and increase to the displayed finite
value, hence converge absolutely. The product has already fixed all
harmonic data. Scalar invariance follows directly from the total
boundary mass one.
\end{proof}

\subsection{Two different tails and the exclusion margin}

For \(R\ge0\), let \(n=\lfloor R\rfloor\), \(q_0=\ee^{-a}\), and
\(S=|X|+13+n\). The explicit two-boundary estimate is
\begin{align}\label{eq:boundary-tail}
 &\int_{|x-X|>R}K(x-X)
  \{|\log|H_\lambda(x)||+|\log|H_\lambda(x+8\ii)||\}\dd x
 \nonumber\\
 &\hspace{1cm}\le
 \frac{C_*}{2}\ee^{-an}S\log^2(3S)
 \frac{1+4q_0+q_0^2}{(1-q_0)^4}.
\end{align}
To derive it, cover both sides by unit shells and use
\(K(t)\le\ee^{-a|t|}/8\) and Lemma~\ref{lem:unit-log} on both
boundaries. The remaining sum is bounded by
\(S\log^2(3S)\sum_{j\ge0}q_0^j(1+j)^3\).
The generating function of this cubic moment is the rational factor
in \eqref{eq:boundary-tail}.

For the sharper zero tail we use the additional canonical unit-count
input from the parent analysis \cite{SmoothLocalizationProof}:
after a fixed enlargement of \(A\),
\[
 \#\{j:u\le x_j(\lambda)\le u+1\}\le A\log(2+|u|)
\]
uniformly for \(0\le\lambda\le1/2\).
This inherited input includes the endpoint \(\lambda=0\); it is
separate from the global Jensen count proved above. Since
\begin{equation}\label{eq:G-decay}
 0<G(t,y)\le K_Gy\ee^{-a|t|},\qquad
 K_G=\frac{2a}{\cos^2a},\quad 0<y\le1,
\end{equation}
unit shells give
\begin{equation}\label{eq:zero-tail}
 \sum_{\substack{y_j>0\\|x_j-X|>R}}G(X-x_j,y_j)
 \le2AK_G\ee^{-an}
 \left\{\frac{\log(3+|X|)+\log(n+1)}{1-q_0}
       +\frac{q_0}{(n+1)(1-q_0)^2}\right\}.
\end{equation}
Here \(\atanh u\le u/(1-u^2)\), \(\sin(ay)\le ay\), and
\(\operatorname{sech}(at)\le2\ee^{-a|t|}\) prove
\eqref{eq:G-decay}. For the shell sum use
\(\log(3+|X|+k)\le\log(3+|X|)+\log(k+1)\) and
\(\log(n+j+1)\le\log(n+1)+j/(n+1)\).
The boundary and zero tails pay different truncations. In particular
\(J_\lambda(X)\ll_A\log(3+|X|)\), but that coarse bound is insufficient
for the following consumer.

\begin{corollary}[A single nonreal-zero margin]\label{cor:margin}
If \(d\le y_j\le1\) and \(|x_j-X|\le r\), then
\begin{equation}\label{eq:margin}
 J_\lambda(X)\ge m(d,r):=\frac{\sin(a)d}{\cosh(ar)}.
\end{equation}
An independently acquired strict upper bound below \(m(d,r)\)
excludes this specified zero class.
\end{corollary}
\begin{proof}
Use positivity of every term, \(\atanh u\ge u\), and the concavity
bound \(\sin(ay)\ge y\sin a\) for \(0\le y\le1\).
\end{proof}

An all-real slice has \(J=0\); any upper zero makes \(J>0\) at every
center. Accordingly universal time-zero vanishing is an RH-equivalent
reformulation, not an additional estimate. Uniform tails on compact
center intervals and continuity of finite root multisets show that
\(J\) is continuous through collisions. A pair approaching a double
real zero contributes a quantity linear in its upper height; when
that height has square-root time scale, so does the Green signal.
Continuity alone does not imply Lipschitz continuity or monotonicity.

